% TOP_PROBLEM: {"id": 30006830, "title": "Sharp Upper Bound in Yau's Conjecture on Nodal Sets of Laplace Eigenfunctions", "category_id": 9, "category": {"id": 9, "name": "pde", "display_name": "Partial Differential Equations", "description": "PDEs and their applications in physics and geometry.", "slug": "pde", "order_index": 9, "created_at": "2026-07-31T15:26:25.671Z"}, "status": "open"}
% FIRST_POSTED: 2026-09-27
% ATTEMPT_STATUS: unresolved
% !TeX program = lualatex
\documentclass[11pt]{article}
\usepackage[margin=25mm]{geometry}
\usepackage{fontspec}
\setmainfont{Latin Modern Roman}
\usepackage{amsmath,amssymb,mathtools}
\usepackage{unicode-math}
\setmathfont{Latin Modern Math}
\usepackage{xurl,hyperref}
\hypersetup{colorlinks=true,urlcolor=blue}
\setcounter{secnumdepth}{0}
\setlength{\parindent}{0pt}
\setlength{\parskip}{5pt}
\setlength{\emergencystretch}{3em}
\begin{document}
\section*{\texorpdfstring{Sharp Upper Bound in Yau's Conjecture on Nodal Sets of Laplace Eigenfunctions}{Sharp Upper Bound in Yau's Conjecture on Nodal Sets of Laplace Eigenfunctions}}\label{30006830}
\subsection{Definitions and mathematical statement}
For the usual connected-manifold formulation, let $\phi\ne0$ satisfy
$-\Delta_g\phi=\lambda\phi$ on a smooth closed Riemannian $n$-manifold, $n\ge2$, and define its nodal set $Z(\phi)=\{x:\phi(x)=0\}$. The target upper estimate is
\[
 \mathcal H^{n-1}(Z(\phi))\le C(M,g)\sqrt\lambda
\]
for every positive eigenvalue and eigenfunction. Hausdorff measure measures the hypersurface-size of the zeros, and the constant is independent of the eigenfunction and eigenvalue. \textbf{Scope qualification:} connectedness is implicit in the standard statement but absent from the catalogue wording. On a disconnected manifold an eigenfunction can vanish identically on an entire component, making the literal bound false; use connected $M$ or explicitly restrict to components where the eigenfunction is nonzero. No analyticity of $g$ is assumed.
\par\smallskip\textit{Formulation and concept references:} \href{https://archive.intlpress.com/site/pub/files/_fulltext/journals/cdm/2018/2018/0001/CDM-2018-2018-0001-a004.pdf}{[S1]}.

\subsection{Short English statement}
On a smooth connected closed manifold, is every eigenfunction's nodal hypersurface measure bounded by a constant times the square root of its eigenvalue?
\subsection{Research attempt}
\paragraph{Scope and source check.}
This unresolved attempt by GPT-6 Astra Ultra was prepared on 27 September
2026 using Fourier calculations, explicit spherical harmonics and the
primary review [S1]. The connectedness qualification in the statement
is retained. The smooth-metric upper problem is distinct from the known
sharp analytic-metric case and from the sharp lower estimate. Results
for Euclidean Dirichlet domains or bi-Laplace equations also have
different hypotheses and are not substituted for this closed-manifold
problem. The attempt proves the sharp order on the standard flat torus
and tests the tempting strategy of a uniform bound in every wavelength
ball.

It is enough to discuss real eigenfunctions. If a complex eigenfunction
is allowed, one of its real and imaginary parts is a nonzero real
eigenfunction with the same eigenvalue, and the complex zero set is
contained in that real zero set. An upper bound for the latter suffices.
All examples and nodal hypersurface calculations below are real-valued.

\paragraph{A complete flat-torus estimate by one-dimensional zero counting.}
Work on $\mathbb T^n=(\mathbb R/2\pi\mathbb Z)^n$ with its standard
flat metric. Fourier expansion of $-\Delta\phi=\lambda\phi$ gives
\[
                 \phi(x)=\sum_{|k|^2=\lambda}a_ke^{ik\cdot x}.
\]
There are finitely many such lattice vectors. Set
$R=\lfloor\sqrt\lambda\rfloor$. On a coordinate circle, keeping
all but $x_i$ fixed, the restriction is a trigonometric polynomial
of degree at most $R$. Unless it vanishes identically on that circle,
it has at most $2R$ distinct zeros: multiply its expression by
$e^{iRx_i}$ and regard it as a polynomial of degree at most $2R$
in $e^{ix_i}$.

For each coordinate direction, the exceptional parameter values giving
an identically zero restriction have $(n-1)$-dimensional measure zero.
Indeed all the coefficient trigonometric polynomials in the other
variables would have to vanish there, and at least one of those
coefficients is nonzero because $\phi\ne0$. The zero set of a
nonzero trigonometric polynomial has measure zero. That last fact can
be proved inductively by the same one-variable zero count and Fubini,
treating identically zero slices in the next lower dimension.

We use the standard nodal-structure fact for nonzero analytic Laplace
eigenfunctions, reviewed in [S1]: the regular nodal hypersurface is
smooth, and its singular part has zero $(n-1)$-dimensional Hausdorff
measure. Let $\nu$ be a unit normal on its regular part. The coarea
formula for projection omitting coordinate $i$ gives
\[
 \int_{Z_{\rm reg}}|\nu_i|\,d\mathcal H^{n-1}
   =\int_{\mathbb T^{n-1}}\#\bigl(Z_{\rm reg}\cap
                  \text{the corresponding coordinate circle}\bigr)
                        \,dx_{\widehat i}
   \le2R(2\pi)^{n-1}.
\]
Exceptional slices do not contribute to this integral. Since a unit
vector satisfies $\sum_i|\nu_i|\ge1$, summing the inequalities yields
\begin{equation}\label{a277:torus}
 \mathcal H^{n-1}(Z(\phi))
                 \le2n(2\pi)^{n-1}\sqrt\lambda.
\end{equation}
Thus the conjectured exponent is proved for every eigenfunction on
this flat torus, including arbitrary linear combinations within a
high-multiplicity eigenspace. No claim that such combinations have
only the zeros of their individual Fourier modes is made.

\paragraph{Exact examples check the scale and constants.}
For $\phi(x)=\sin(mx_1)$, there are $2m$ parallel nodal subtori,
each of measure $(2\pi)^{n-1}$, and $\lambda=m^2$. Hence
\[
 \mathcal H^{n-1}(Z(\phi))=2(2\pi)^{n-1}\sqrt\lambda.
\]
The square-root order cannot be improved even within this family.
For another check on the two-dimensional torus,
$\phi(x,y)=\sin(mx)\sin(\ell y)$ has eigenvalue $m^2+\ell^2$.
Its nodal set is a union of $2m$ vertical and $2\ell$ horizontal
circles. Their finitely many intersections have zero length, so
\[
             \mathcal H^1(Z)=4\pi(m+\ell)
                    \le4\pi\sqrt2\sqrt{m^2+\ell^2}.
\]
This confirms that adding crossing families changes the constant but
not the exponent. Counting each grid intersection as another curve
segment would double-count a measure-zero part.

The normalization matters. Using period one instead of $2\pi$ changes
both the eigenvalue formula and the geometric measure. The estimate
above fixes both conventions before counting zeros, so no frequency
factor $2\pi$ is lost when comparing with $\sqrt\lambda$.

\paragraph{A proposed wavelength-ball shortcut fails on the sphere.}
Suppose one could cover a general manifold by $O(r^{-n})$ balls of
radius $r=\lambda^{-1/2}$ and bound nodal measure in every such ball
by $Cr^{n-1}$, with $C$ depending only on the metric. Summation would
immediately give $C/r=C\sqrt\lambda$. Unfortunately the proposed
local bound is false, including on the analytic round sphere where
the desired global bound is true.

On $S^2\subset\mathbb R^3$, take
\[
                 \phi_m=\operatorname{Re}(x_1+ix_2)^m.
\]
The homogeneous polynomial is harmonic in $\mathbb R^3$, so its
restriction has eigenvalue $m(m+1)$. Its zeros are $m$ distinct great
circles through the north and south poles. Distinct circles intersect
only at the poles, a length-zero set. Thus the total nodal length is
$2\pi m$, of the desired global order.

But a geodesic ball of radius $r<\pi/2$ centered at the north pole
contains two radial arcs of length $r$ from each circle. Its nodal
length is exactly $2mr$. At wavelength radius
$r=[m(m+1)]^{-1/2}$ this remains comparable to two, whereas the
proposed uniform local bound $Cr$ tends to zero. The quotient of the
actual local length by $r$ is $2m$, unbounded. High vanishing order
at one point obstructs a uniform per-ball constant even though it
does not violate the global conjecture.

This example also explains why a compactness argument over normalized
wavelength-scale functions cannot assume bounded vanishing order.
The metric is fixed throughout, while the eigenvalue and the local
multiplicity increase. Normalizing an $L^2$ norm does not remove the
many meridian sheets near the pole.

\paragraph{A more realistic conditional summation criterion.}
Let $r=c\lambda^{-1/2}$ be below a fixed coordinate radius and choose
a covering by $O(r^{-n})$ balls. Suppose one has numbers $N_j\ge1$
measuring local oscillation or doubling, together with two estimates
\[
 \mathcal H^{n-1}(Z\cap B_j)\le C N_j r^{n-1},
 \qquad \sum_jN_j\le C' r^{-n}.
\]
Then summing proves
$\mathcal H^{n-1}(Z)\le CC' r^{-1}=O(\sqrt\lambda)$.
This deduction is exact. Neither displayed hypothesis is asserted here
for all smooth metrics. The first is a linear local nodal estimate;
the second is a bound on the average complexity across the cover,
not a uniform bound on each $N_j$.

The spherical example is compatible with exceptional large $N_j$ near
a small set of poles, but not with their uniform boundedness. A worst-case
global frequency bound at each ball, multiplied by the full number
of balls, loses precisely the distributional information needed for
the sharp exponent. Known polynomial estimates in frequency do not
turn into linear estimates by renaming their constants. A successful
implementation must account quantitatively for how high-complexity
balls are distributed.

\paragraph{Harmonic lifting is useful but does not supply that estimate.}
Define on the product manifold $M\times\mathbb R$
\[
                   u(x,t)=e^{\sqrt\lambda t}\phi(x).
\]
For the product metric,
$\Delta_{M\times\mathbb R}u=e^{\sqrt\lambda t}
(\Delta_M\phi+\lambda\phi)=0$.
Its zero set is $Z(\phi)\times\mathbb R$. Thus harmonic-function
estimates can be applied to the lifted problem, with explicit
exponential growth in the new direction. On a product interval of
length $L$, the regular zero-set measure is
$L\mathcal H^{n-1}(Z(\phi))$, with the singular part negligible in
the standard nodal framework.

The lift does not make the metric analytic when $g$ is only smooth.
Nor does it guarantee a linear relationship between a doubling index
and nodal measure. It is a correct change of equation that transfers
the central quantitative difficulty. The Fourier proof of
\eqref{a277:torus} used a finite frequency band on every coordinate
circle; restrictions of eigenfunctions to geodesics for a general
smooth metric need not be trigonometric polynomials of bounded degree.
Applying their one-dimensional polynomial zero count would therefore
be unjustified.

\paragraph{Boundary cases and verification.}
Connectedness prevents a nonzero eigenfunction from vanishing identically
on an entire separate component, the literal obstruction already noted
in the catalogue. Positive eigenvalues exclude the constant-function
case. Finitely many low eigenvalues can be absorbed into a metric-
dependent constant when using the wavelength cover; this does not
address the high-frequency estimate. The proof and tests separately
track hypersurface measure, isolated crossing points, eigenvalue
normalization and singular nodal points.

Exact finite checks verify the torus grid formulas, the eigenvalue and
local length of the spherical example, and the powers of $r$ in the
conditional sum. The analytic nodal-structure and coarea inputs are
stated explicitly. The strongest result is the all-eigenfunction flat-
torus estimate, plus a concrete disproof of a tempting uniform local
shortcut. The first remaining gap is a sharp quantitative local-to-global
mechanism for arbitrary smooth metrics, such as the two-part summation
package above. The catalogue's smooth upper-bound assertion remains
\textbf{UNRESOLVED}.

\subsection{Sources}
\begin{itemize}
\item[S1]A. Logunov, E. Malinnikova, 'Review of Yau's conjecture on zero sets of Laplace eigenfunctions,' Current Developments in Mathematics 2018, pp. 179-212, International Press (2019).
\url{https://archive.intlpress.com/site/pub/files/_fulltext/journals/cdm/2018/2018/0001/CDM-2018-2018-0001-a004.pdf}.
\textit{Formulation source.}
\end{itemize}
\end{document}
